% GENERATED FILE. Edit canonical lessons, then run npm run generate:books.
% canonical-source-hash: fnv1a64:04cb964594f47865
% canonical-lessons: KA-S167-letter-dha, KA-S168-letter-bha, KA-S169-letter-pha, KA-S170-letter-kha, KA-S171-letter-gha, KA-S172-letter-ddha, KA-S173-letter-ttha, KA-R80-the-breath

\chapter{The Breath}
\label{ch:ka-usiru}
\hlchaptermodality{80}

\begin{chapteropening}
\textbf{By the end of this chapter:} \emph{I can recognise the seven Kannada consonants that carry the breath, and say which plain letter each one is built on.}

From cold: pair each of the seven breathy consonants with the plain letter it is built on, and say why the romanization's ph is not the f of phone.
\end{chapteropening}

\section[\texorpdfstring{\kn{ಧ} (dha)}{ಧ (dha)}]{\kn{ಧ} — the breath that opens the thank-you word}

\label{lesson:KA-S167-letter-dha}



\begin{quote}
Say thank you in Kannada. It begins on today's character, and it has done since the first chapter.
\end{quote}

\subsection*{Script you'll notice: \kn{ಧ}}
\textbf{\kn{ಧ}} — \emph{dha}.

It is a \textbf{consonant}, and in this script a consonant is never bare: it comes with an \emph{a} already in it. So it is not \emph{dh}, it is \textbf{dha}.

\textbf{\kn{ಧ}} is \textbf{\kn{ದ}} with a puff of air after it. Hold a hand in front of your mouth: say \kn{ದ} and the air stays put; say \kn{ಧ} and it comes out.

\textbf{This is the last big idea left in this script, and it accounts for seven of the characters you have never been shown.} Kannada gives the breathy partner of a consonant its own shape rather than leaving you to guess:

\noindent
\begin{tabularx}{\linewidth}{@{}>{\raggedright\arraybackslash}X>{\raggedright\arraybackslash}X@{}}
\toprule
\textbf{no breath} & \textbf{breath} \\
\midrule
\kn{ದ} & \kn{ಧ} \\
\bottomrule
\end{tabularx}

The six that follow are the same move on six other letters.

You already say these:

\begin{itemize}
\raggedright
  \item \textbf{\kn{ಧನ್ಯವಾದ}} \emph{dhanyavāda} — thank you
  \item \textbf{\kn{ಮಧ್ಯಾಹ್ನ}} \emph{madhyāhna} — noon
  \item \textbf{\kn{ಶುಭ} \kn{ಮಧ್ಯಾಹ್ನ}} \emph{śubha madhyāhna} — good afternoon
\end{itemize}

\subsection*{Writing: \kn{ಧ}}
\hlblockfigure{\detokenize{figures/KA-S167-letter-dha-filmstrip.pdf}}{How \kn{ಧ} is written, stroke by stroke}

\begin{itemize}
\raggedright
  \item \textbf{1.} start at the upper left and go down the left side into the left lobe, rising into the middle point
  \item \textbf{2.} without lifting, drop from the point around the right lobe and climb the right side
  \item \textbf{3.} without lifting, close the bowl leftward along the top
  \item \textbf{4.} lift, then draw the top bar from left to right
  \item \textbf{5.} without lifting, curl up into the hook
  \item \textbf{6.} lift, then draw the tail downward
\end{itemize}

\textbf{Pen lifts: 2.} This is \kn{ದ} with a tail: the body, then the curled bar, then the tail below the point, each after a lift.

\begin{quote}
This is one attested teaching order and not a national standard — Kannada handwriting is taught with school-to-school variation. Source: Gopala Krishna A, 'Kannada-alphabet-dhha.gif', 42 frames, Wikimedia Commons, 25 May 2016, CC BY-SA 4.0.
\end{quote}

\begin{quote}
The animation is filed as 'dhha': that series spells \kn{ದ} as 'dha'.
\end{quote}

\subsection*{Your turn}
\begin{itemize}
\raggedright
  \item \emph{Look:} at these words, and find \kn{ಧ} in each
\end{itemize}

\begin{quote}
\kn{ಧನ್ಯವಾದ}  ·  \kn{ಮಧ್ಯಾಹ್ನ}
\end{quote}

\begin{itemize}
\raggedright
  \item \emph{Trace it:} \kn{ಧ} three times, saying \emph{dha} as you finish each one
  \item \emph{Look:} back at any page of this book and find \kn{ಧ} once more
\end{itemize}

\begin{tcolorbox}[breakable,colback=teal!4,colframe=teal!35!black,title={Before you move on}]
Which character is this — \kn{ಧ}? What sound does it carry? (\emph{\textbf{dha}}.) What does it add to \kn{ದ}? (\textbf{The breath after it}.)
\end{tcolorbox}
\section[\texorpdfstring{\kn{ಭ} (bha)}{ಭ (bha)}]{\kn{ಭ} — the breath inside \kn{ಶುಭ}}

\label{lesson:KA-S168-letter-bha}



\begin{quote}
Before the new one: \kn{ಧ} — what does it add to \kn{ದ}? Today's character does the same thing to \kn{ಬ}.
\end{quote}

\subsection*{Script you'll notice: \kn{ಭ}}
\textbf{\kn{ಭ}} — \emph{bha}.

It is a \textbf{consonant}, and in this script a consonant is never bare: it comes with an \emph{a} already in it. So it is not \emph{bh}, it is \textbf{bha}.

\noindent
\begin{tabularx}{\linewidth}{@{}>{\raggedright\arraybackslash}X>{\raggedright\arraybackslash}X@{}}
\toprule
\textbf{no breath} & \textbf{breath} \\
\midrule
\kn{ಬ} & \kn{ಭ} \\
\bottomrule
\end{tabularx}

\textbf{\kn{ಶುಭ}} \emph{śubha} carries it, so every greeting in this book ends on it — and you learned \kn{ಶ} for the front of that word one page ago. Today the back of it arrives.

You already say these:

\begin{itemize}
\raggedright
  \item \textbf{\kn{ಶುಭ} \kn{ರಾತ್ರಿ}} \emph{śubha rātri} — good night
  \item \textbf{\kn{ಶುಭೋದಯ}} \emph{śubhōdaya} — good morning
  \item \textbf{\kn{ಭುಜ}} \emph{bhuja} — the shoulder
\end{itemize}

\subsection*{Writing: \kn{ಭ}}
\hlblockfigure{\detokenize{figures/KA-S168-letter-bha-filmstrip.pdf}}{How \kn{ಭ} is written, stroke by stroke}

\begin{itemize}
\raggedright
  \item \textbf{1.} start at the curled tip inside the head and loop up over it clockwise
  \item \textbf{2.} without lifting, slant down to the left and round the left lobe, rising into the middle point
  \item \textbf{3.} without lifting, drop round the right lobe and climb the right side to the bar
  \item \textbf{4.} without lifting, run back left along the bar
  \item \textbf{5.} without lifting, draw the bar from left to right
  \item \textbf{6.} without lifting, curl up into the hook
  \item \textbf{7.} lift, then draw the tail downward
\end{itemize}

\textbf{Pen lifts: 1.} The body is drawn as \kn{ಬ}'s is; then the pen runs back along the short bar and draws it into the hook. It lifts once, for the tail. The animation also lifts before the bar; most people copying \kn{ಭ} in the Omniglot study used two strokes.

\begin{quote}
This is one attested teaching order and not a national standard — Kannada handwriting is taught with school-to-school variation. Source: Gopala Krishna A, 'Kannada-alphabet-bha.gif', 33 frames, Wikimedia Commons, 25 May 2016, CC BY-SA 4.0.
\end{quote}

\subsection*{Your turn}
\begin{itemize}
\raggedright
  \item \emph{Look:} at these words, and find \kn{ಭ} in each
\end{itemize}

\begin{quote}
\kn{ಶುಭ} \kn{ರಾತ್ರಿ}  ·  \kn{ಭುಜ}
\end{quote}

\begin{itemize}
\raggedright
  \item \emph{Trace it:} \kn{ಭ} three times, saying \emph{bha} as you finish each one
  \item \emph{Look:} back at any page of this book and find \kn{ಭ} once more
\end{itemize}

\begin{tcolorbox}[breakable,colback=teal!4,colframe=teal!35!black,title={Before you move on}]
Which character is this — \kn{ಭ}? What sound does it carry? (\emph{\textbf{bha}}.) Which letter is it the breathy partner of? (\textbf{\kn{ಬ}}.)
\end{tcolorbox}
\section[\texorpdfstring{\kn{ಫ} (pha)}{ಫ (pha)}]{\kn{ಫ} — and the two English letters that mislead}

\label{lesson:KA-S169-letter-pha}



\begin{quote}
Before the new one: \kn{ಭ} — which letter is it the breathy partner of? Today's does the same to \kn{ಪ}.
\end{quote}

\subsection*{Script you'll notice: \kn{ಫ}}
\textbf{\kn{ಫ}} — \emph{pha}.

It is a \textbf{consonant}, and in this script a consonant is never bare: it comes with an \emph{a} already in it. So it is not \emph{ph}, it is \textbf{pha}.

\noindent
\begin{tabularx}{\linewidth}{@{}>{\raggedright\arraybackslash}X>{\raggedright\arraybackslash}X@{}}
\toprule
\textbf{no breath} & \textbf{breath} \\
\midrule
\kn{ಪ} & \kn{ಫ} \\
\bottomrule
\end{tabularx}

\textbf{Read the romanization carefully here.} \emph{ph} is one sound with air behind it, not the \emph{f} those two English letters make in \emph{phone}. Two letters on the romanized line, one character on the Kannada line, one puff in the mouth.

The word below is the clearest case, because English lent it and Kannada spells it with this letter.

You already say these:

\begin{itemize}
\raggedright
  \item \textbf{\kn{ಕಾಫಿ}} \emph{kāphi} — coffee
  \item \textbf{\kn{ನಿಮಗೆ} \kn{ಕಾಫಿ} \kn{ಇಷ್ಟವಾ}} \emph{nimage kāphi iṣṭavā} — do you like coffee?
\end{itemize}

\subsection*{Writing: \kn{ಫ}}
\hlblockfigure{\detokenize{figures/KA-S169-letter-pha-filmstrip.pdf}}{How \kn{ಫ} is written, stroke by stroke}

\begin{itemize}
\raggedright
  \item \textbf{1.} start at the inner tip of the curl and wind up and round it to the left
  \item \textbf{2.} without lifting, run along the base and rise into the middle point
  \item \textbf{3.} without lifting, drop round the right lobe and climb the right side to its tip
  \item \textbf{4.} lift, then set the dot in the middle
  \item \textbf{5.} lift, then draw the top bar from left to right
  \item \textbf{6.} without lifting, curl up into the hook
  \item \textbf{7.} lift, then draw the tail downward
\end{itemize}

\textbf{Pen lifts: 3.} This is \kn{ಪ} with a tail: the body, the dot, the curled bar, and last the tail, each after a lift.

\begin{quote}
This is one attested teaching order and not a national standard — Kannada handwriting is taught with school-to-school variation. Source: Gopala Krishna A, 'Kannada-alphabet-pha.gif', 43 frames, Wikimedia Commons, 25 May 2016, CC BY-SA 4.0.
\end{quote}

\subsection*{Your turn}
\begin{itemize}
\raggedright
  \item \emph{Look:} at these words, and find \kn{ಫ} in each
\end{itemize}

\begin{quote}
\kn{ಕಾಫಿ}
\end{quote}

\begin{itemize}
\raggedright
  \item \emph{Trace it:} \kn{ಫ} three times, saying \emph{pha} as you finish each one
  \item \emph{Look:} back at any page of this book and find \kn{ಫ} once more
\end{itemize}

\begin{tcolorbox}[breakable,colback=teal!4,colframe=teal!35!black,title={Before you move on}]
Which character is this — \kn{ಫ}? What sound does it carry? (\emph{\textbf{pha}}.) Is it the \emph{f} of \emph{phone}? (\textbf{No} — a \emph{p} with breath after it.)
\end{tcolorbox}
\section[\texorpdfstring{\kn{ಖ} (kha)}{ಖ (kha)}]{\kn{ಖ} — the breath on the first letter you ever wrote}

\label{lesson:KA-S170-letter-kha}



\begin{quote}
Before the new one: \kn{ಫ} — which letter is it built on? Today's is built on \kn{ಕ}, the consonant this book drew first.
\end{quote}

\subsection*{Script you'll notice: \kn{ಖ}}
\textbf{\kn{ಖ}} — \emph{kha}.

It is a \textbf{consonant}, and in this script a consonant is never bare: it comes with an \emph{a} already in it. So it is not \emph{kh}, it is \textbf{kha}.

\noindent
\begin{tabularx}{\linewidth}{@{}>{\raggedright\arraybackslash}X>{\raggedright\arraybackslash}X@{}}
\toprule
\textbf{no breath} & \textbf{breath} \\
\midrule
\kn{ಕ} & \kn{ಖ} \\
\bottomrule
\end{tabularx}

This one is rarer than the three before it, and the word that needs it is one you already use to agree with somebody.

You already say these:

\begin{itemize}
\raggedright
  \item \textbf{\kn{ಖಂಡಿತ}} \emph{khaṇḍita} — certainly
\end{itemize}

\subsection*{Writing: \kn{ಖ}}
\hlblockfigure{\detokenize{figures/KA-S170-letter-kha-filmstrip.pdf}}{How \kn{ಖ} is written, stroke by stroke}

\begin{itemize}
\raggedright
  \item \textbf{1.} start at the inner tip of the curl and wind out of it, up and over to the right
  \item \textbf{2.} without lifting, come down through the waist to the base
  \item \textbf{3.} without lifting, round the lower loop and cross back through the waist
  \item \textbf{4.} without lifting, run along the base and climb the right side, curving over to its tip
\end{itemize}

\textbf{Pen lifts: 0.} The whole letter is one run, from the curl through the lower loop to the tip of the long right side. There is no curled bar on top.

\begin{quote}
This is one attested teaching order and not a national standard — Kannada handwriting is taught with school-to-school variation. Source: Gopala Krishna A, 'Kannada-alphabet-kha.gif', 37 frames, Wikimedia Commons, 25 May 2016, CC BY-SA 4.0.
\end{quote}

\subsection*{Your turn}
\begin{itemize}
\raggedright
  \item \emph{Look:} at these words, and find \kn{ಖ} in each
\end{itemize}

\begin{quote}
\kn{ಖಂಡಿತ}
\end{quote}

\begin{itemize}
\raggedright
  \item \emph{Trace it:} \kn{ಖ} three times, saying \emph{kha} as you finish each one
  \item \emph{Look:} back at any page of this book and find \kn{ಖ} once more
\end{itemize}

\begin{tcolorbox}[breakable,colback=teal!4,colframe=teal!35!black,title={Before you move on}]
Which character is this — \kn{ಖ}? What sound does it carry? (\emph{\textbf{kha}}.) Which word from the agreeing chapter opens on it? (\textbf{\kn{ಖಂಡಿತ}}.)
\end{tcolorbox}
\section[\texorpdfstring{\kn{ಘ} (gha)}{ಘ (gha)}]{\kn{ಘ} — rare, and worth a page anyway}

\label{lesson:KA-S171-letter-gha}



\begin{quote}
Before the new one: \kn{ಖ} — which letter is it built on? Today's is built on \kn{ಗ}.
\end{quote}

\subsection*{Script you'll notice: \kn{ಘ}}
\textbf{\kn{ಘ}} — \emph{gha}.

It is a \textbf{consonant}, and in this script a consonant is never bare: it comes with an \emph{a} already in it. So it is not \emph{gh}, it is \textbf{gha}.

\noindent
\begin{tabularx}{\linewidth}{@{}>{\raggedright\arraybackslash}X>{\raggedright\arraybackslash}X@{}}
\toprule
\textbf{no breath} & \textbf{breath} \\
\midrule
\kn{ಗ} & \kn{ಘ} \\
\bottomrule
\end{tabularx}

\textbf{This character appears in one word in this whole book}, and that is the reason to teach it rather than a reason to skip it. The word is a month name, and a reader who meets an unfamiliar shape inside a list they already know by heart cannot tell whether the word or the printing is at fault.

You already say these:

\begin{itemize}
\raggedright
  \item \textbf{\kn{ಮಾಘ}} \emph{māgha} — Magha, the eleventh traditional month
\end{itemize}

\subsection*{Writing: \kn{ಘ}}
\hlblockfigure{\detokenize{figures/KA-S171-letter-gha-filmstrip.pdf}}{How \kn{ಘ} is written, stroke by stroke}

\begin{itemize}
\raggedright
  \item \textbf{1.} start at the inner tip of the curl and wind up and round it to the left
  \item \textbf{2.} without lifting, run along the base and rise into the middle point
  \item \textbf{3.} without lifting, drop round the right lobe and climb the right side
  \item \textbf{4.} without lifting, carry on up the middle arm to its tip
  \item \textbf{5.} without lifting, come back down the middle arm to its foot
  \item \textbf{6.} without lifting, round the right arm and curl in at its top
  \item \textbf{7.} lift, then draw the top bar from left to right
  \item \textbf{8.} without lifting, curl up into the hook
  \item \textbf{9.} lift, then set the dot in the middle
  \item \textbf{10.} lift, then draw the tail downward
\end{itemize}

\textbf{Pen lifts: 3.} After the middle arm the pen comes back down it and rounds the right arm; then the curled bar, the dot and the tail each come after a lift. The animation also lifts between the arms; most people copying \kn{ಘ} in the Omniglot study used four strokes.

\begin{quote}
This is one attested teaching order and not a national standard — Kannada handwriting is taught with school-to-school variation. Source: Gopala Krishna A, 'Kannada-alphabet-gha.gif', 61 frames, Wikimedia Commons, 25 May 2016, CC BY-SA 4.0.
\end{quote}

\subsection*{Your turn}
\begin{itemize}
\raggedright
  \item \emph{Look:} at these words, and find \kn{ಘ} in each
\end{itemize}

\begin{quote}
\kn{ಮಾಘ}
\end{quote}

\begin{itemize}
\raggedright
  \item \emph{Trace it:} \kn{ಘ} three times, saying \emph{gha} as you finish each one
  \item \emph{Look:} back at any page of this book and find \kn{ಘ} once more
\end{itemize}

\begin{tcolorbox}[breakable,colback=teal!4,colframe=teal!35!black,title={Before you move on}]
Which character is this — \kn{ಘ}? What sound does it carry? (\emph{\textbf{gha}}.) Which month carries it? (\textbf{\kn{ಮಾಘ}}.)
\end{tcolorbox}
\section[\texorpdfstring{\kn{ಢ} (ḍha)}{ಢ (ḍha)}]{\kn{ಢ} — the curled d, with the breath}

\label{lesson:KA-S172-letter-ddha}



\begin{quote}
Before the new one: \kn{ಘ} — which letter is it built on? Today's is built on \kn{ಡ}, the curled d.
\end{quote}

\subsection*{Script you'll notice: \kn{ಢ}}
\textbf{\kn{ಢ}} — \emph{ḍha}.

It is a \textbf{consonant}, and in this script a consonant is never bare: it comes with an \emph{a} already in it. So it is not \emph{ḍh}, it is \textbf{ḍha}.

\noindent
\begin{tabularx}{\linewidth}{@{}>{\raggedright\arraybackslash}X>{\raggedright\arraybackslash}X@{}}
\toprule
\textbf{no breath} & \textbf{breath} \\
\midrule
\kn{ಡ} & \kn{ಢ} \\
\bottomrule
\end{tabularx}

Two things at once: the tongue is \textbf{curled back}, the way it is for \kn{ಣ} and \kn{ಷ}, and the breath comes \textbf{after}. Like \kn{ಘ}, it is in one word in this book, and that word is a month.

You already say these:

\begin{itemize}
\raggedright
  \item \textbf{\kn{ಆಷಾಢ}} \emph{āṣāḍha} — Ashadha, the fourth traditional month
\end{itemize}

\subsection*{Writing: \kn{ಢ}}
\hlblockfigure{\detokenize{figures/KA-S172-letter-ddha-filmstrip.pdf}}{How \kn{ಢ} is written, stroke by stroke}

\begin{itemize}
\raggedright
  \item \textbf{1.} start at the upper left and go down the left side into the left lobe, rising into the middle point
  \item \textbf{2.} without lifting, drop from the point around the right lobe and climb the right side
  \item \textbf{3.} without lifting, curl left into the small inner loop and go round it counterclockwise
  \item \textbf{4.} without lifting, rise out of the loop and close the bowl leftward along the top
  \item \textbf{5.} lift, then draw the top bar from left to right
  \item \textbf{6.} without lifting, curl up into the hook
  \item \textbf{7.} lift, then draw the tail downward
\end{itemize}

\textbf{Pen lifts: 2.} This is \kn{ಡ} with a tail: the body, then the curled bar, then the tail below the point, each after a lift.

\begin{quote}
This is one attested teaching order and not a national standard — Kannada handwriting is taught with school-to-school variation. Source: Gopala Krishna A, 'Kannada-alphabet-dda.gif', 52 frames, Wikimedia Commons, 25 May 2016, CC BY-SA 4.0.
\end{quote}

\begin{quote}
The animation is filed as 'dda': that series spells \kn{ಡ} as 'da'.
\end{quote}

\subsection*{Your turn}
\begin{itemize}
\raggedright
  \item \emph{Look:} at these words, and find \kn{ಢ} in each
\end{itemize}

\begin{quote}
\kn{ಆಷಾಢ}
\end{quote}

\begin{itemize}
\raggedright
  \item \emph{Trace it:} \kn{ಢ} three times, saying \emph{ḍha} as you finish each one
  \item \emph{Look:} back at any page of this book and find \kn{ಢ} once more
\end{itemize}

\begin{tcolorbox}[breakable,colback=teal!4,colframe=teal!35!black,title={Before you move on}]
Which character is this — \kn{ಢ}? What sound does it carry? (\emph{\textbf{ḍha}}.) Two things make it different from \kn{ದ} — name them. (\textbf{The tongue curls back, and the breath follows}.)
\end{tcolorbox}
\section[\texorpdfstring{\kn{ಠ} (ṭha)}{ಠ (ṭha)}]{\kn{ಠ} — the last one}

\label{lesson:KA-S173-letter-ttha}



\begin{quote}
Before the new one: \kn{ಢ} — what two things make it what it is? Today's is the same pair of moves on \kn{ಟ}.
\end{quote}

\subsection*{Script you'll notice: \kn{ಠ}}
\textbf{\kn{ಠ}} — \emph{ṭha}.

It is a \textbf{consonant}, and in this script a consonant is never bare: it comes with an \emph{a} already in it. So it is not \emph{ṭh}, it is \textbf{ṭha}.

\noindent
\begin{tabularx}{\linewidth}{@{}>{\raggedright\arraybackslash}X>{\raggedright\arraybackslash}X@{}}
\toprule
\textbf{no breath} & \textbf{breath} \\
\midrule
\kn{ಟ} & \kn{ಠ} \\
\bottomrule
\end{tabularx}

\textbf{This is the last character in this book that nothing had taught.} It appears three times, all of them in the same month name. From here on, every Kannada character these pages print is one somebody has handed you.

You already say these:

\begin{itemize}
\raggedright
  \item \textbf{\kn{ಜ್ಯೇಷ್ಠ}} \emph{jyēṣṭha} — Jyeshtha, the third traditional month
\end{itemize}

\subsection*{Writing: \kn{ಠ}}
\hlblockfigure{\detokenize{figures/KA-S173-letter-ttha-filmstrip.pdf}}{How \kn{ಠ} is written, stroke by stroke}

\begin{itemize}
\raggedright
  \item \textbf{1.} start at the upper left of the bowl and go down the left side and round the base
  \item \textbf{2.} without lifting, climb the right side
  \item \textbf{3.} without lifting, close the bowl leftward along the top
  \item \textbf{4.} lift, then draw the top bar from left to right
  \item \textbf{5.} without lifting, curl up into the hook
  \item \textbf{6.} lift, then set the dot in the middle
\end{itemize}

\textbf{Pen lifts: 2.} The round bowl is one run, closing back along the top as in \kn{ದ}; then come the curled bar and last the dot, each after a lift.

\begin{quote}
This is one attested teaching order and not a national standard — Kannada handwriting is taught with school-to-school variation. Source: Gopala Krishna A, 'Kannada-alphabet-tta.gif', 41 frames, Wikimedia Commons, 25 May 2016, CC BY-SA 4.0.
\end{quote}

\begin{quote}
The animation is filed as 'tta': that series spells \kn{ಟ} as 'ta' and \kn{ಥ} as 'thha'.
\end{quote}

\subsection*{Your turn}
\begin{itemize}
\raggedright
  \item \emph{Look:} at these words, and find \kn{ಠ} in each
\end{itemize}

\begin{quote}
\kn{ಜ್ಯೇಷ್ಠ}
\end{quote}

\begin{itemize}
\raggedright
  \item \emph{Trace it:} \kn{ಠ} three times, saying \emph{ṭha} as you finish each one
  \item \emph{Look:} back at any page of this book and find \kn{ಠ} once more
\end{itemize}

\begin{tcolorbox}[breakable,colback=teal!4,colframe=teal!35!black,title={Before you move on}]
Which character is this — \kn{ಠ}? What sound does it carry? (\emph{\textbf{ṭha}}.) Which letter is it the breathy partner of? (\textbf{\kn{ಟ}}.)
\end{tcolorbox}
\section[\texorpdfstring{(\kn{ಖ} \kn{ಘ} \kn{ಠ} \kn{ಢ} \kn{ಧ} \kn{ಫ} \kn{ಭ})}{(ಖ ಘ ಠ ಢ ಧ ಫ ಭ)}]{(\kn{ಖ} \kn{ಘ} \kn{ಠ} \kn{ಢ} \kn{ಧ} \kn{ಫ} \kn{ಭ}) — seven characters, one move}

\label{lesson:KA-R80-the-breath}



\begin{quote}
Close the lessons before this one. Say thank you, and say good morning. Name the breathy character inside each.
\end{quote}

\subsection*{Script: seven shapes, one move}
Seven characters, and \textbf{they are one move done seven times}: take a consonant you already read and let the breath out after it. Hold a hand in front of your mouth and the difference is not subtle.

\noindent
\begin{tabularx}{\linewidth}{@{}>{\raggedright\arraybackslash}X>{\raggedright\arraybackslash}X@{}}
\toprule
\textbf{no breath} & \textbf{breath} \\
\midrule
\kn{ಕ} \emph{ka} & \kn{ಖ} \emph{kha} \\
\kn{ಗ} \emph{ga} & \kn{ಘ} \emph{gha} \\
\kn{ಟ} \emph{ṭa} & \kn{ಠ} \emph{ṭha} \\
\kn{ಡ} \emph{ḍa} & \kn{ಢ} \emph{ḍha} \\
\kn{ದ} \emph{da} & \kn{ಧ} \emph{dha} \\
\kn{ಪ} \emph{pa} & \kn{ಫ} \emph{pha} \\
\kn{ಬ} \emph{ba} & \kn{ಭ} \emph{bha} \\
\bottomrule
\end{tabularx}

Read the right-hand column against the left and the shapes stop being seven new things to memorise. \textbf{English does not write this difference at all}, which is why the romanization has to spell it with an added \emph{h} — and why \emph{ph} on that line is a \emph{p} with breath, never the \emph{f} of \emph{phone}.

Two of the seven are in one word each, and both words are months: \textbf{\kn{ಮಾಘ}} and \textbf{\kn{ಆಷಾಢ}}. \textbf{\kn{ಠ}} is in \textbf{\kn{ಜ್ಯೇಷ್ಠ}}, three times in the whole book, and it was the last character these pages printed that nothing had taught.

\subsection*{Your turn}
\begin{itemize}
\raggedright
  \item \emph{Look:} at these words and find the breathy character in each
\end{itemize}

\begin{quote}
\kn{ಧನ್ಯವಾದ}  ·  \kn{ಶುಭೋದಯ}  ·  \kn{ಕಾಫಿ}  ·  \kn{ಜ್ಯೇಷ್ಠ}
\end{quote}

Say these aloud:

\begin{itemize}
\raggedright
  \item \kn{ದ} then \kn{ಧ}, and \kn{ಬ} then \kn{ಭ}, with a hand in front of your mouth
  \item which of the seven you have seen most often on these pages
\end{itemize}

\begin{tcolorbox}[breakable,colback=teal!4,colframe=teal!35!black,title={Before you move on}]
Which letter is \kn{ದ} with the breath? (\textbf{\kn{ಧ}}.) Which is \kn{ಬ} with the breath? (\textbf{\kn{ಭ}}.) Is \kn{ಫ} the \emph{f} of \emph{phone}? (\textbf{No} — a \emph{p} with breath after it.) Which of the seven was the last character in this book with nothing teaching it? (\textbf{\kn{ಠ}}.)
\end{tcolorbox}
